\documentclass[11pt]{article}
\usepackage[margin=1in]{geometry}
\usepackage{amsmath,amssymb,amsthm}
\usepackage{enumitem}

\newtheorem{definition}{Definition}
\newtheorem{example}{Example}
\newtheorem{proposition}{Proposition}

\title{Module 2: Basic Probability}
\author{Mathematical Models of Infectious Diseases}
\date{}

\begin{document}
\maketitle

\section{Sample spaces and events}

\begin{definition}
An \emph{experiment} is a process with uncertain outcome. The
\emph{sample space} $\Omega$ is the set of all possible outcomes. An
\emph{event} is a subset $A\subseteq\Omega$.
\end{definition}

\begin{example}
A susceptible individual has one contact with an infectious person.
\[
\Omega=\{\text{infected},\ \text{not infected}\}.
\]
$A=\{\text{infected}\}$ is an event.
\end{example}

\section{Set operations on events}

\begin{itemize}
  \item \textbf{Union} $A\cup B$: ``$A$ or $B$ occurs.''
  \item \textbf{Intersection} $A\cap B$: ``$A$ and $B$ occur.''
  \item \textbf{Complement} $A^c$: ``$A$ does not occur.''
  \item \textbf{Disjoint}: $A\cap B=\emptyset$ ($A$ and $B$ cannot occur together).
\end{itemize}

\begin{example}
$A=\{$infected with strain 1$\}$, $B=\{$infected with strain 2$\}$.
\begin{itemize}[nosep]
  \item $A\cup B$: infected with at least one strain.
  \item $A\cap B$: co-infected with both strains.
  \item $A^c$: not infected with strain 1.
\end{itemize}
\end{example}

\section{The probability function}

\begin{definition}[Axioms of probability]
$P$ assigns to each event $A\subseteq\Omega$ a number $P(A)\in[0,1]$ such that
\begin{itemize}[nosep]
  \item $P(\Omega)=1$;
  \item for disjoint events $A,B$, $\ P(A\cup B)=P(A)+P(B)$.
\end{itemize}
\end{definition}

Interpretation: $P(A)$ is the long-run fraction of times $A$ occurs (e.g. $P(\text{infection})=0.3$ means 30\% of exposed contacts result in infection).

\section{Probability of the complement}
A very useful identity is $\Pr(A^c) = 1-P(A)$.

%% \begin{proposition}[Addition law]
%% $P(A\cup B) = P(A)+P(B)-P(A\cap B)$.
%% \end{proposition}

\section{Independence}

\begin{definition}
Events $A$ and $B$ are \emph{independent} if
\[
P(A \cap B) = P(A)\,P(B).
\]
\end{definition}

Interpretation: knowing that some of the events occurred gives no information about whether the others occurred.

\begin{example}
A susceptible individual has $n$ contacts. Each contact independently transmits infection with probability $c$. Let $A_i=\{\text{contact } i \text{ does not transmit infection}\}$, so $P(A_i)=1-c$. Avoiding infection on all $n$ contacts is the event $A_1\cap\cdots\cap A_n$, so by independence
\[
P(\text{avoid infection}) = \prod_{i=1}^n P(A_i) = (1-c)^n.
\]
\end{example}

\section{Random variables}

\begin{definition}
A \emph{random variable} $X$ is a function $X:\Omega\to\mathbb{R}$ assigning a number to each outcome.
\end{definition}

\begin{example}
$T=$ time until recovery.
\end{example}

\begin{definition}
The \emph{density} (pdf) of $X$ is a function $f\ge 0$ with $\displaystyle\int_{-\infty}^\infty f(x)\,dx=1$ such that
\[
P(a\le X\le b)=\int_a^b f(x)\,dx.
\]
The \emph{expectation} is $E[X]=\int x f(x)\,dx$.
\end{definition}

\section{The exponential random variable}

Used to model waiting times: time to recovery, time until next contact, incubation period.

\begin{definition}
$T$ has an \emph{exponential distribution with rate $\lambda>0$}, written $T\sim\text{Exp}(\lambda)$, if
\[
f(t)=\lambda e^{-\lambda t}, \qquad t\ge 0.
\]
\end{definition}

The \emph{survival function} is the probability of not yet having recovered by time $t$:
\[
S(t) = P(T>t) = \int_t^\infty \lambda e^{-\lambda s}\,ds = e^{-\lambda t}.
\]
Mean: $E[T]=1/\lambda$.

Interpretation: if $T$ is the infectious period and $T\sim\text{Exp}(\gamma)$, then $\gamma$ is the recovery \emph{rate} and $1/\gamma$ is the mean infectious period.

\section{The memoryless property}

\begin{definition}
$T$ is \emph{memoryless} if
\[
P(T\ge t+s \mid T\ge t) = P(T\ge s) \qquad \text{for all } s,t\ge 0.
\]
\end{definition}

Interpretation: given that an individual is still infectious at time $t$, the probability they remain infectious for $s$ more days does not depend on $t$. Time already spent infected gives no information about time remaining --- the recovery rate is constant, with no ``aging'' or ``wearing out.''

\begin{proposition}
If $T\sim\text{Exp}(\lambda)$, then $T$ is memoryless.
\end{proposition}

\begin{proof}
Since $t+s\ge t$, the event $\{T\ge t+s\}\subseteq\{T\ge t\}$, so $\{T\ge t+s\}\cap\{T\ge t\}=\{T\ge t+s\}$. Hence
\[
P(T\ge t+s\mid T\ge t) = \frac{P(T\ge t+s)}{P(T\ge t)} = \frac{e^{-\lambda(t+s)}}{e^{-\lambda t}} = e^{-\lambda s} = P(T\ge s). \qquad\qed
\]
\end{proof}

Consequence: the exponential is the \emph{only} continuous distribution with this property. This is why, in the simplest infectious disease models, an exponentially distributed infectious period corresponds to a constant per-capita recovery rate $\gamma$, independent of how long an individual has already been infected.

\end{document}
